\section{Zadanie 13 (Hubert Jastrzębski)}
Rzucamy $n$ kul w sposób jednostajny i niezależny do $n$ urn. Niech $X_i$ będzie indykatorem na zdarzenie: $i$-ta urna jest pusta. Niech $Y_1, \dots, Y_n$ będą niezależnymi zmiennymi wskaźnikowymi z prawdopodobieństwem sukcesu $p = \left(1 - \frac{1}{n}\right)^n$. Niech $X = \sum_{i=1}^n X_i$ oraz $Y = \sum_{i=1}^n Y_i$.

\begin{enumerate}
    \item[(i)] Pokaż, że dla każdego $k \in \mathbb{N}_1$, zachodzi $E(X_1 X_2 \dots X_k) \le E(Y_1 Y_2 \dots Y_k)$.
    \item[(ii)] Pokaż, że dla każdego $t \in \mathbb{N}$, zachodzi $E(e^{tX}) \le E(e^{tY})$.
    \item[(iii)] Wyznacz ograniczenie Czernowa dla $P(X \ge (1 + \delta) E(X))$.
\end{enumerate}

\subsection*{Rozwiązanie}

\textbf{(i)}

Po prostu liczymy obydwie strony:
\begin{enumerate}
    \item $E(X_1 X_2 \dots X_k) = \left(\frac{n-k}{n}\right)^n = \left(1 - \frac{k}{n}\right)^n$, ponieważ aby wszystkie z $k$ urn były puste, każda z $n$ kul musi trafić do jednej z pozostałych $n-k$ urn.
    \item $E(Y_1 Y_2 \dots Y_k) = p^k =  \left(1 - \frac{1}{n}\right)^{nk}$, ponieważ $Y_i$ są niezależne i no tak jest.
\end{enumerate}
Każda strona jest dodatnia, więc sobie podnosimy do potęgi $\frac{1}{n}$:
\[1-\frac{k}{n} \le \left(1 - \frac{1}{n}\right)^{k}\]

Nierówność Bernoulliego mówi, że \(1 + kx \le (1+x)^k, \text{ dla } x \ge -1 \land k \in \mathbb{N}\).

Podstawiając \(x = -\frac{1}{n}\) (\(n \ge 1 \implies x = -\frac{1}{n} \ge -1\)), dostajemy:
\[1 - \frac{k}{n} \le \left(1 - \frac{1}{n}\right)^{k}\]
czyli dokładnie to co chcieliśmy udowodnić.

\textbf{(ii)}
\[
e^{t X_i} = \left\{
\begin{array}{ll}
e^t, & \text{dla } X_i = 1 \\
1, & \text{dla } X_i = 0
\end{array}
\right. = e^t X_i + (1 - X_i) = 1 + (e^t - 1) X_i
\]

Zatem:
\[
e^{tX} = \prod_{i=1}^n e^{t X_i} = \prod_{i=1}^n \left[1 + (e^t - 1) X_i\right]
= \sum_{S \subseteq [n]} (e^t - 1)^{|S|} \prod_{i \in S} X_i
= \sum_{k=0}^n (e^t - 1)^k \sum_{|S| = k} \prod_{i \in S} X_i
\]
To przejście z sumą to po prostu wymnażanie wielomianów.

Bierzemy teraz wartość oczekiwaną:
\[
E(e^{tX}) = \sum_{k=0}^n (e^t - 1)^k \sum_{|S| = k} E\left(\prod_{i \in S} X_i\right)
\]

Podobnie wyprowadzamy prawą część:
\[
E(e^{tY}) = \sum_{k'=0}^n (e^t - 1)^{k'} \sum_{|S'| = k'} E\left(\prod_{i \in S'} Y_i\right)
\]

Z (i) wiemy, że dla każdego \(m = |S| = |S'|\) zachodzi:
\[E\left(\prod_{i \in S} X_i\right) \le E\left(\prod_{i \in S'} Y_i\right)\]

A że wszystko jest nieujemne, to mamy, że dla każdego \(t \in \mathbb{N}\) zachodzi:
\[
\sum_{k=0}^n (e^t - 1)^k \sum_{|S| = k} E\left(\prod_{i \in S} X_i\right)
\le
\sum_{k'=0}^n (e^t - 1)^{k'} \sum_{|S'| = k'} E\left(\prod_{i \in S'} Y_i\right)
\]
\[E(e^{tX}) \le E(e^{tY})\]

Dodatkowo, co nam się przyda w (iii), możemy zauważyć, że tak na prawdę to zachodzi
dla każdego \(t \ge 0\), bo \(\forall_{t \ge 0} \left( e^{t} - 1 \ge 0 \right) \implies \forall_{t \ge 0, k \ge 0} \left( (e^{t} - 1)^{k} \ge 0\right)\)

\textbf{(iii)}

Na początku jeszcze rozpiszmy sobie dokładniej \(E(e^{tY})\) (korzystając z (ii)):
\[
E(e^{tY}) = \sum_{k=0}^n (e^t - 1)^{k} \sum_{|S| = k} E\left(\prod_{i \in S} Y_i\right)
\]
zmienne \(Y_i\) są niezależne, więc:
\[
E(e^{tY}) = \sum_{k=0}^n (e^t - 1)^{k} \sum_{|S| = k} \prod_{i \in S} E(Y_i)
\]
\[
E(e^{tY}) = \sum_{k=0}^n (e^t - 1)^{k} \sum_{|S| = k} p^{k} = \sum_{k=0}^n (e^t - 1)^{k} \binom{n}{k} p^{k}
\]
\[
E(e^{tY}) = \sum_{k=0}^n \binom{n}{k} \cdot \left(p(e^t - 1)\right)^{k} \cdot 1^{n-k} = \left(1 + p(e^t - 1)\right)^{n}
\]

Teraz przejdźmy do ograniczenia Czernowa:
\[
\forall_{t \ge 0} P(X \ge (1 + \delta) E(X)) \le \frac{E(e^{tX})}{e^{t(1 + \delta) E(X)}} \le \frac{E(e^{tY})}{e^{t(1 + \delta) E(X)}} = \frac{\left(1 + p(e^t - 1)\right)^{n}}{e^{t(1 + \delta) E(X)}}
\]
Pierwszą nierówność mamy z Czernowa, drugą z (ii), a równość przed chwilą pokazaliśmy.

Pozostaje nam wyliczyć \(E(X)\):
\[
E(X) = \sum_{i=1}^n E(X_i) = \sum_{i=1}^n P(X_i = 1) = \sum_{i=1}^n \left(1 - \frac{1}{n}\right)^n = n \left(1 - \frac{1}{n}\right)^n = n p
\]

Tak więc ostatecznie mamy:
\[
\frac{\left(1 + p(e^t - 1)\right)^{n}}{e^{t(1 + \delta) np}} = \left(\frac{1 + p(e^t - 1)}{e^{tp(1 + \delta)}}\right)^n
\]

Ale chyba teraz niestety trzeba jeszcze znaleźć minimum tej funkcji względem \(t\) \dots

Skorzystamy z nierówności \(1 + x \le e^x\) dla \(x = p(e^t - 1)\):
\[
\left(\frac{1 + p(e^t - 1)}{e^{tp(1 + \delta)}}\right)^n \le \left(\frac{e^{p(e^t - 1)}}{e^{tp(1 + \delta)}}\right)^n = \left(e^{p(e^t - 1 - t(1 + \delta))}\right)^n = e^{np(e^t - 1 - t(1 + \delta))}
\]

\[
f(t) = e^{np(e^t - 1 - t(1 + \delta))}
\]

Chcemy znaleźć minimum tej funkcji. Liczymy pochodną:
\[
f'(t) = np e^{t} - np(1 + \delta) = 0 \implies e^{t} = 1 + \delta \implies t = \ln(1 + \delta)
\]

Druga pochodna:
\[
f''(t) = np e^{t} > 0 \quad \forall_{t}
\]

Czyli znaleźliśmy minimum. Podstawiamy z powrotem do \(f(t)\):
\begin{align*}
f(\ln(1 + \delta)) &= e^{np((1 + \delta) - 1 - \ln(1 + \delta)(1 + \delta))} \\
&= e^{np(\delta - (1 + \delta) \ln(1 + \delta))} \\
&= \left( e^{\delta} \cdot \left(e^{\ln(1+\delta)}\right)^{-(1 + \delta)} \right)^{np} \\
&= \left( e^{\delta} \cdot \left(1+\delta\right)^{-(1 + \delta)} \right)^{np} \\
&= \left( \frac{e^{\delta}}{\left(1+\delta\right)^{(1 + \delta)}} \right)^{np}
\end{align*}

\subsection*{Odpowiedź}
\[
\boxed{\text{(iii) \( \left( \frac{e^{\delta}}{\left(1+\delta\right)^{(1 + \delta)}} \right)^{np} \)}}
\]
